\section{Contoh Terhitung}
\label{sec:contoh-ecc}

Contoh-contoh pada bagian ini disusun mengikuti alur bab, dari aritmetika
medan sampai protokol yang lengkap. Contoh pertama bekerja di dalam
$GF(2^8)$, tempat seluruh aturan medan berhingga terlihat paling jelas.
Contoh kedua dan ketiga menjalankan penjumlahan serta penggandaan titik pada
kurva, masing-masing di atas medan berhingga dan di atas bilangan riil.
Tiga contoh terakhir merangkai seluruh perkakas itu menjadi protokol
pertukaran kunci, enkripsi, dan pengkodean pesan. Setiap angka di bawah ini
telah diverifikasi dengan program pada bagian praktikum di akhir bab ini.

\begin{example}[Aritmetika polinomial di dalam medan $GF(2^8)$]
\label{ex:gf2m}
Setiap \textit{byte} dipandang sebagai polinom berderajat paling tinggi tujuh,
dengan setiap bit menjadi koefisien. Dengan kesepakatan itu,
\[
\code{0D} = x^3 + x^2 + 1,
\qquad
\code{06} = x^2 + x .
\]

\textbf{(a) Penjumlahan.} Penjumlahan di dalam $GF(2)$ dilakukan modulo dua,
sehingga $1 + 1 = 0$; koefisien yang sama pada kedua polinom saling
menghapus. Karena itu penjumlahan dua polinom tidak lain adalah operasi XOR
bit demi bit. Untuk polinom di atas,
\[
(x^3 + x^2 + 1) \oplus (x^2 + x) = x^3 + (x^2 \oplus x^2) + x + 1
= x^3 + x + 1 = \code{0B} .
\]
Suku $x^2$ lenyap karena muncul pada kedua polinom. Dengan cara yang sama,
\[
\code{57} \oplus \code{83} = \code{D4} ,
\]
sebab dalam biner $\code{57} = 01010111$, $\code{83} = 10000011$, dan
XOR keduanya adalah $11010100 = \code{D4}$.

\textbf{(b) Perkalian.} Perkalian dilakukan dalam dua tahap: kalikan dahulu
seperti polinom biasa, lalu reduksi hasilnya dengan polinom reduksi
$x^8 + x^4 + x^3 + x + 1$ (yaitu $\code{11B}$, polinom yang dipakai AES).
Tulis kedua faktor dalam bentuk polinom,
\[
\code{57} = x^6 + x^4 + x^2 + x + 1,
\qquad
\code{83} = x^7 + x + 1 .
\]
Karena $\code{83}$ memiliki bit pada posisi $7$, $1$, dan $0$, perkaliannya
adalah XOR dari tiga salinan $\code{57}$ yang digeser sejauh posisi bit
tersebut. Tabel~\ref{tab:gf2m-kali} memperlihatkan langkah itu.

\begin{table}[htbp]
    \centering
    \caption{Perkalian dua polinom di dalam $GF(2^8)$: menghasilkan polinom
    mentah berderajat $13$, lalu direduksi menjadi berderajat $7$}
    \label{tab:gf2m-kali}
    \begin{tabular}{@{}l l l@{}}
    \toprule
    \textbf{Langkah} & \textbf{Biner} & \textbf{Bentuk polinom} \\
    \midrule
    $\code{57} \ll 0$        & $01010111$          & $x^6 + x^4 + x^2 + x + 1$ \\
    \addlinespace
    $\code{57} \ll 1$        & $10101110$          & $x^7 + x^5 + x^3 + x^2 + x$ \\
    \addlinespace
    $\code{57} \ll 7$        & $10101110000000$    & $x^{13} + x^{11} + x^9 + x^8 + x^7$ \\
    \midrule
    XOR ketiganya            & $10101101111001$    & $x^{13} + x^{11} + x^9 + x^8 + x^6
                             + x^5 + x^4 + x^3 + 1$ \\
    \midrule
    $\oplus\,(\code{11B} \ll 5)$ & $0000100000011001$ & $x^{11} + x^4 + x^3 + 1$ \\
    \addlinespace
    $\oplus\,(\code{11B} \ll 3)$ & $11000001$         & $x^7 + x^6 + 1 = \code{C1}$ \\
    \bottomrule
    \end{tabular}
\end{table}

Baris keempat Tabel~\ref{tab:gf2m-kali} adalah hasil kali mentah, yang
berderajat $13$ sehingga masih di luar medan. Dua baris terakhir
memperlihatkannya direduksi: selama derajatnya paling sedikit delapan,
polinom reduksi digeser ke posisi bit terdepan lalu di-XOR-kan. Hasil
akhirnya adalah $x^7 + x^6 + 1$, yaitu $\code{C1}$. Jadi
\[
\code{57} \cdot \code{83} = \code{C1} .
\]

\textbf{(c) Pemeriksaan pada medan yang lebih kecil.} Dengan cara yang sama,
di dalam $GF(2^4)$ yang memakai polinom reduksi $x^4 + x + 1$,
\[
(x^3 + x^2 + 1)(x^2 + x) = x^5 + x^3 + x^2 + x ,
\]
sebab koefisien $x^4$ dan $x^3$ yang muncul dua kali saling menghapus. Hasil
kali itu berderajat lima, sehingga direduksi satu langkah:
\[
(x^5 + x^3 + x^2 + x) \oplus \bigl(x \cdot (x^4 + x + 1)\bigr)
= (x^5 + x^3 + x^2 + x) \oplus (x^5 + x^2 + x) = x^3 .
\]
Sisanya berderajat tiga, jadi hasil akhirnya adalah $x^3 = \code{1000}$.

\par\vspace{2pt}{\footnotesize\itshape Catatan: dua angka pada berkas sumber
perlu dikoreksi. Pertama, berkas sumber menuliskan hasil kali mentah di atas
sebagai $\code{10110}$, padahal bentuk binernya adalah $\code{101110}$; dua
bit terdepan $\code{10}$ hilang. Kedua, berkas sumber menyebut $x^2 + 1$
tidak dapat direduksi, padahal di dalam $GF(2)$ polinom itu sama dengan
$(x+1)^2$, sebab $(x+1)^2 = x^2 + 2x + 1 = x^2 + 1$. Karena itu berkas
sumber juga salah ketika memfaktorkan $x^5 + x^2 + x + 1$ menjadi
$(x^3 + x^2 + 1)(x^2 + 1)$; hasil kali kedua faktor itu adalah
$x^5 + x^4 + x^3 + 1$, bukan $x^5 + x^2 + x + 1$. Faktorisasi yang benar
adalah $x^5 + x^2 + x + 1 = (x^2 + 1)(x^3 + x + 1)$. Kedua koreksi itu
diperiksa oleh program pada bagian praktikum.\par}
\end{example}

\begin{example}[Penjumlahan titik pada kurva di atas $GF(11)$]
\label{ex:gf11}
Ambil kurva
\[
y^2 \equiv x^3 + x + 6 \pmod{11},
\]
yang ordonya $13$ dan telah dienumerasi pada Subbagian~\ref{subsec:enumerasi-gf-p}.
Dua di antara titiknya adalah $P(2, 4)$ dan $Q(5, 9)$; keduanya memenuhi
kurva, sebab $4^2 = 16 \equiv 5$ dan $2^3 + 2 + 6 = 16 \equiv 5$ untuk $P$,
serta $9^2 = 81 \equiv 4$ dan $5^3 + 5 + 6 = 136 \equiv 4$ untuk $Q$.

\textbf{(a) Penjumlahan $P + Q$.} Dengan $a = 1$ dan $p = 11$, gradiennya
menurut Persamaan~\eqref{eq:gradien-mod-p} adalah
\[
m \equiv (9 - 4)(5 - 2)^{-1} \equiv 5 \cdot 3^{-1} \pmod{11} .
\]
Karena $3 \cdot 4 = 12 \equiv 1$, invers dari $3$ adalah $4$, sehingga
$m \equiv 5 \cdot 4 = 20 \equiv 9$. Selanjutnya
\[
x_r \equiv m^2 - x_p - x_q \equiv 81 - 2 - 5 = 74 \equiv 8 \pmod{11},
\]
\[
y_r \equiv m(x_p - x_r) - y_p \equiv 9(2 - 8) - 4 = -58 \equiv 8 \pmod{11},
\]
sebab $-58 + 6 \cdot 11 = 8$. Jadi $P + Q = (8, 8)$, dan pemeriksaannya:
$8^2 = 64 \equiv 9$ sedangkan $8^3 + 8 + 6 = 526 = 47 \cdot 11 + 9$. Cocok.

\textbf{(b) Penggandaan $2P$.} Rumusnya sama, hanya gradiennya diganti
menurut Persamaan~\eqref{eq:gradien-penggandaan}:
\[
m \equiv (3 \cdot 2^2 + 1)(2 \cdot 4)^{-1} \equiv 13 \cdot 8^{-1}
\pmod{11} .
\]
Karena $8 \cdot 7 = 56 \equiv 1$, invers dari $8$ adalah $7$, sehingga
$m \equiv 13 \cdot 7 = 91 \equiv 3$. Koordinat hasilnya
\[
x_r \equiv m^2 - 2x_p \equiv 9 - 4 = 5, \qquad
y_r \equiv m(x_p - x_r) - y_p \equiv 3(2 - 5) - 4 = -13 \equiv 9 ,
\]
jadi $2P = (5, 9)$.

Perhatikan hasil yang menarik: $2P$ ternyata sama dengan $Q$. Artinya $Q$
bukan titik yang bebas, melainkan kelipatan dari $P$, dan karena itu
$P + Q = 3P = (8, 8)$. Kebetulan semacam ini lazim pada kurva kecil dan
menjadi alasan mengapa pada kurva yang dipakai sungguhan orderya harus
diperiksa lebih dahulu. Bila titik basis berorde kecil, hubungan seperti ini
memberi petunjuk kepada penyerang. Tabel pelelaran $kP$ untuk
$k = 1, 2, \dots, 13$ pada kurva ini, yang berakhir di $13P = \mathcal{O}$,
diperlihatkan pada Tabel~\ref{tab:pelelaran-gf11}.
\end{example}

\begin{example}[Penjumlahan dan penggandaan titik di atas bilangan riil]
\label{ex:real}
Kurva di atas bilangan riil dipakai untuk memahami geometrinya, bukan untuk
kriptografi. Dua kurva berikut memperlihatkan kedua operasi itu.

\textbf{(a) Kurva $y^2 = x^3 - 3x + 3$.} Di sini $a = -3$ dan $b = 3$,
sehingga $4(-3)^3 + 27(3)^2 = -108 + 243 = 135 \neq 0$ menurut
Persamaan~\eqref{eq:diskriminan}: kurva itu licin. Titik $P(1, 1)$ memenuhi
kurva karena $1^2 = 1 - 3 + 3 = 1$, dan titik $Q(-2, 1)$ juga karena
$1^2 = -8 + 6 + 3 = 1$.

Untuk $P + Q$ dengan $P \neq Q$,
\[
m = \frac{1 - 1}{-2 - 1} = \frac{0}{-3} = 0 ,
\]
sehingga garis yang melalui $P$ dan $Q$ mendatar. Akibatnya
\[
x_r = 0 - 1 - (-2) = 1, \qquad y_r = 0 \cdot (1 - 1) - 1 = -1 ,
\]
dan $P + Q = (1, -1)$. Titik itu benar pada kurva sebab
$(-1)^2 = 1 = 1 - 3 + 3$.

Untuk $2P$,
\[
m = \frac{3(1)^2 - 3}{2(1)} = \frac{0}{2} = 0 ,
\]
sehingga garis singgungnya pun mendatar. Maka $x_r = 0 - 2 = -2$ dan
$y_r = 0 \cdot (1 + 2) - 1 = -1$, jadi $2P = (-2, -1)$. Periksa:
$(-1)^2 = (-2)^3 - 3(-2) + 3 = -8 + 6 + 3 = 1$. Cocok.

\textbf{(b) Kurva $y^2 = x^3 + 2x + 4$.} Di sini $a = 2$ dan $b = 4$.
Titik $P(2, 4)$ memenuhi kurva karena $4^2 = 8 + 4 + 4 = 16$, dan
$Q(0, 2)$ juga karena $2^2 = 0 + 0 + 4 = 4$. Gradien garis $PQ$ adalah
\[
m = \frac{2 - 4}{0 - 2} = \frac{-2}{-2} = 1 ,
\]
sehingga $x_r = 1 - 2 - 0 = -1$ dan $y_r = 1 \cdot (2 + 1) - 4 = -1$.
Jadi $P + Q = (-1, -1)$, dan titik itu benar pada kurva sebab
$(-1)^2 = 1 = -1 - 2 + 4$.

Penggandaannya menghasilkan pecahan:
\[
m = \frac{3(2)^2 + 2}{2(4)} = \frac{14}{8} = \frac{7}{4},
\]
\[
x_r = \left(\frac{7}{4}\right)^2 - 2(2) = \frac{49}{16} - 4 = -\frac{15}{16},
\]
\[
y_r = \frac{7}{4}\left(2 + \frac{15}{16}\right) - 4
    = \frac{7 \cdot 47}{64} - 4 = \frac{329 - 256}{64} = \frac{73}{64} .
\]
Jadi $2P = \left(-\frac{15}{16}, \frac{73}{64}\right)$, atau dalam pecahan
desimal $(-0{,}9375;\ 1{,}140625)$. Bandingkan dengan kurva di atas
$GF(11)$: di sana koordinat hasilnya tetap bilangan bulat, sedangkan di sini
koordinatnya langsung menjadi pecahan. Itulah sebabnya kriptografi memakai
kurva di atas medan berhingga.
\end{example}

\begin{example}[Pertukaran kunci ECDH]
\label{ex:ecdh}
\textbf{(a) Pada kurva di atas $GF(5)$.} Ambil kurva
\[
y^2 \equiv x^3 + 2x + 1 \pmod{5},
\]
yang memiliki $6$ titik berhingga ditambah $\mathcal{O}$, sehingga orderya
$7$. Titik basis yang disepakati adalah $B(0, 1)$, yang memenuhi kurva sebab
$1^2 = 0 + 0 + 1$. Alice memilih kunci privat $a = 2$ dan Bob $b = 3$.
Kunci publik masing-masing adalah
\[
P_A = 2B = (1, 3),
\qquad
P_B = 3B = (3, 3).
\]
Setiap pihak menghitung kunci bersama dari kunci publik pihak lain:
\[
K = a \cdot P_B = 2(3, 3) = (0, 4),
\qquad
K = b \cdot P_A = 3(1, 3) = (0, 4) .
\]
Keduanya memperoleh titik $(0, 4)$ --- perhatikan bahwa titik itu sama dengan
$-B$, sebab $B$ berorde $7$ sehingga $6B = -B$.

\textbf{(b) Pada kurva di atas $GF(23)$.} Kurva $y^2 \equiv x^3 + x + 1
\pmod{23}$ dengan titik basis $B(3, 10)$ memberi angka yang lebih besar.
Dengan $a = 2$ dan $b = 3$,
\[
P_A = 2B = (7, 12),
\qquad
P_B = 3B = (19, 5),
\]
\[
K = 2P_B = (12, 4),
\qquad
K = 3P_A = (12, 4) .
\]
Penyerang yang menyadap pertukaran itu hanya melihat kurva, titik basis
$B(3,10)$, dan kedua kunci publik $(7,12)$ serta $(19,5)$. Untuk menemukan
$K$ ia harus memulihkan $a$ atau $b$ dari $P_A = aB$, yaitu memecahkan
persoalan logaritma diskrit kurva eliptik pada Persamaan~\eqref{eq:ecdlp}.
Pada kurva mainan ini penyerangan itu sepele --- program pada bagian
praktikum melakukannya dengan penelusuran langkah bayi-langkah raksasa ---
tetapi pada kurva 256 bit usaha yang sama menuntut sekitar $2^{128}$ langkah.
\end{example}

\begin{example}[Enkripsi dan dekripsi ECEG]
\label{ex:eceg}
Enkripsi kurva eliptik memakai kurva yang sama dengan contoh sebelumnya, yaitu
$y^2 \equiv x^3 + x + 1 \pmod{23}$ dengan titik basis $B(3, 10)$. Kunci
privat penerima adalah $b = 3$, sehingga kunci publiknya
\[
P_B = 3B = (19, 5) .
\]
Pesan yang hendak dikirim telah dikodekan lebih dahulu menjadi sebuah titik
pada kurva (caranya diperlihatkan pada contoh berikutnya); di sini misalkan
$P_M = (7, 12)$.

Pengirim memilih kunci sesaat $k = 5$ --- bilangan acak yang baru untuk
setiap pesan --- lalu menghitung dua titik menurut
Persamaan~\eqref{eq:eceg-cipher}:
\[
C_1 = kB = 5B = (9, 16),
\]
\[
C_2 = P_M + k P_B = (7, 12) + 5(19, 5) = (7, 12) + (1, 16) = (18, 3) .
\]
Jadi cipherteksnya adalah pasangan titik $\bigl((9, 16),\ (18, 3)\bigr)$.

Penerima memegang $b = 3$ dan memulihkan pesannya dengan mengurangi titik
kedua dengan hasil kali kunci privatnya dan titik pertama:
\[
P_M = C_2 - b\,C_1 = (18, 3) - 3(9, 16) .
\]
Karena $C_1 = 5B$, maka $3C_1 = 15B$, dan $15B = (1, 16)$ menurut kenyataan
bahwa $B$ berorde $28$. Lawan titik itu adalah $(1, 7)$, sebab
$-16 \equiv 7 \pmod{23}$, sehingga
\[
P_M = (18, 3) + (1, 7) = (7, 12) .
\]
Plainteksnya pulih tepat seperti semula.

Perhatikan mengapa langkah itu selalu berhasil. Penerima tidak perlu mengetahui
$k$; ia hanya perlu bahwa $kP_B = kbB = bkB = bC_1$, sehingga suku acak pada
$C_2$ terhapus oleh suku yang ia hitung sendiri --- persis seperti penurunan
pada Persamaan~\eqref{eq:eceg-dekripsi}. Sekalipun penyerang menyadap kedua
titik cipherteks, ia tidak dapat melakukan hal yang sama karena tidak memegang
$b$.

Satu hal terakhir: karena $k$ dipilih baru untuk setiap pesan, titik
$C_1 = kB$ juga berubah setiap kali. Pesan $P_M$ yang sama karena itu
menghasilkan cipherteks yang berbeda pada setiap pengiriman. Sifat itu,
yang disebut enkripsi probabilistik, membuat penyerang tidak dapat mengenali
pesan yang berulang --- keunggulan yang tidak dimiliki buku kode maupun
pemetaan karakter langsung pada Subbagian~\ref{subsec:pemetaan-langsung}.
\end{example}

\begin{example}[Pengkodean pesan dengan metode Koblitz]
\label{ex:koblitz}
Agar pesan dapat dienkripsi dengan ECEG, rangkaian bitnya harus diubah
terlebih dahulu menjadi titik pada kurva. Metode Koblitz melakukan itu dengan
memanfaatkan peluang: ruas kanan kurva hanya punya peluang sekitar setengah
untuk menjadi kuadrat modulo $p$ pada nilai $x$ yang dipilih acak.

Ambil kurva
\[
y^2 \equiv x^3 - x + 188 \pmod{751},
\]
yang memiliki $N = 727$ titik. Menurut kesepakatan pengkodean pada
Subbagian~\ref{subsec:koblitz}, huruf \texttt{A} sampai \texttt{Z} bernilai
$10$ sampai $35$, sehingga huruf \texttt{B} bernilai $m = 11$. Ambil
parameter basis $k = 20$, sehingga nilai $x$ pertama yang dicoba adalah
\[
x = mk + 1 = (11)(20) + 1 = 221
\]
menurut Persamaan~\eqref{eq:koblitz-x}. Sulihkan ke dalam kurva:
\[
221^3 - 221 + 188 = 10\,793\,861 - 33 = 10\,793\,828
\equiv 456 \pmod{751} .
\]
Apakah $456$ kuadrat modulo $751$? Pemeriksaan pada program menunjukkan tidak,
sehingga $x = 221$ tidak memberi titik. Pencarian berlanjut seperti pada
Tabel~\ref{tab:koblitz-b}: $x = 222$ memberi $446$ dan $x = 223$ memberi
$266$, keduanya juga bukan kuadrat. Baru pada percobaan keempat berhasil.
Untuk $x = 224$ diperoleh
\[
224^3 - 224 + 188 = 11\,239\,424 - 36 = 11\,239\,388
\equiv 673 \pmod{751} ,
\]
dan $673$ adalah kuadrat modulo $751$: dua akarnya adalah $y = 248$ dan
$y = 503$, sebab $248^2 = 61\,504 = 81 \cdot 751 + 673$ dan
$503^2 = 253\,009 = 336 \cdot 751 + 673$. Jadi huruf \texttt{B} dikodekan
menjadi titik $(224, 248)$ pada kurva tersebut.

Penerima memulihkan pesannya dengan Persamaan~\eqref{eq:koblitz-dekode},
\[
m = \left\lfloor \frac{224 - 1}{20} \right\rfloor
  = \left\lfloor \frac{223}{20} \right\rfloor
  = \lfloor 11{,}15 \rfloor = 11 ,
\]
yang memang nilai untuk huruf \texttt{B}.

Empat percobaan pada contoh ini sejalan dengan dugaan peluangnya. Setiap
percobaan berpeluang sekitar setengah untuk berhasil, sehingga rata-rata
dibutuhkan dua percobaan, dan peluang menuntut lebih dari sepuluh percobaan
lebih kecil dari $10^{-3}$. Agar pemetaan itu tetap satu-satu, kurva harus
memiliki paling sedikit $36k$ titik, sebab setiap nilai $m$ memakai wilayah
$x$ selebar $k$ dan ada $36$ nilai $m$ yang mungkin.
\end{example}
